
\documentclass[11pt]{article}
\usepackage[letterpaper,margin=0.72in]{geometry}
\usepackage[T1]{fontenc}
\usepackage{lmodern}
\usepackage{microtype}
\usepackage{booktabs,tabularx,array}
\usepackage{xcolor}
\usepackage{enumitem}
\usepackage{titlesec}
\usepackage{fancyhdr}
\usepackage{tcolorbox}
\usepackage{hyperref}
\usepackage{siunitx}

\definecolor{navy}{HTML}{16324F}
\definecolor{blue}{HTML}{245B8A}
\definecolor{light}{HTML}{EEF4F8}
\definecolor{accent}{HTML}{B84A3A}
\definecolor{graytext}{HTML}{4D5963}

\hypersetup{colorlinks=true,linkcolor=blue,urlcolor=blue}
\pagestyle{fancy}
\fancyhf{}
\lhead{\small\color{graytext}Treadmill Conditioning}
\rhead{\small\color{graytext}Jordan Suber}
\cfoot{\small\color{graytext}\thepage}
\setlength{\headheight}{14pt}
\setlength{\parindent}{0pt}
\setlength{\parskip}{5pt}
\setlist[itemize]{leftmargin=1.25em,itemsep=2pt,topsep=3pt}
\titleformat{\section}{\Large\bfseries\color{navy}}{}{0pt}{}
\titleformat{\subsection}{\large\bfseries\color{blue}}{}{0pt}{}
\newcolumntype{Y}{>{\raggedright\arraybackslash}X}

\begin{document}

{\color{navy}\Huge\bfseries Hard Treadmill Conditioning}\\[3pt]
{\Large Personalized training handout}\\[7pt]
{\color{graytext}\textbf{Athlete:} Jordan Suber \quad
\textbf{Age:} 31 (32 on November 5) \quad
\textbf{Height:} 6'1" \quad
\textbf{Body mass:} 190 lb}

\begin{tcolorbox}[colback=light,colframe=navy,boxrule=0.8pt,arc=2mm]
\textbf{Programming objective.}
Build high-end aerobic power, speed endurance, and threshold capacity without turning every treadmill session into an all-out test. Intensity is prescribed primarily by \textbf{RPE, breathing, and repeatability}; treadmill speed is individualized.
\end{tcolorbox}

\section{Your Intensity Framework}

For age 31, the Tanaka age-predicted equation ($208-0.7\times age$) gives an estimated maximum heart rate of approximately \textbf{186 bpm}. This is only a population estimate; your actual HRmax may differ materially.

\begin{center}
\begin{tabularx}{\textwidth}{@{}l c c Y@{}}
\toprule
\textbf{Zone / use} & \textbf{\% est. HRmax} & \textbf{Approx. bpm} & \textbf{Practical cue}\\
\midrule
Easy / recovery & 60--70\% & 112--130 & Comfortable; full conversation\\
Steady aerobic & 70--80\% & 130--149 & Controlled; sentences possible\\
Hard / threshold-like & 80--90\% & 149--168 & RPE 7--8; short phrases\\
HIIT peak range & 90--95\% & 168--177 & RPE 8--9; only a few words\\
\bottomrule
\end{tabularx}
\end{center}

\textbf{Important:} Heart rate lags during short intervals. Do not accelerate merely because the monitor has not yet reached the target. For 30--90 second work bouts, RPE, mechanics, and your ability to repeat the effort are more useful controls.

\section{Workout A: 60/120 High-Intensity Intervals}

\textbf{Use:} Primary high-intensity treadmill conditioning session. \textbf{Total time:} about 30--35 minutes.

\subsection{1. Progressive warm-up --- 10 minutes}
\begin{itemize}
\item 0--3 min: brisk walk to easy jog, RPE 2--3.
\item 3--7 min: easy run, RPE 3--4.
\item 7--10 min: gradually increase pace; include 2--3 brief 15--20 sec pickups with easy running between them.
\end{itemize}

\subsection{2. Main set --- start with 6 rounds}
\begin{tcolorbox}[colback=white,colframe=accent,boxrule=1pt,arc=2mm]
\centering
{\Large\bfseries 60 sec HARD \quad + \quad 120 sec EASY}\\[3pt]
Hard interval: RPE 8--9/10 \qquad Recovery: RPE 2--3/10
\end{tcolorbox}

\begin{itemize}
\item Select a speed that is \textbf{fast but mechanically controlled}, not an arbitrary 7.5--10+ mph target.
\item By later rounds, heart rate may approach roughly \textbf{168--177 bpm}; it need not do so during every interval.
\item The final 10--15 seconds should be difficult, but running form should remain stable.
\item Walk or jog on the moving belt during recovery. \textbf{Do not routinely jump onto the side rails.}
\item If round 6 remains crisp and repeatable, progress toward 7 and then 8 rounds over subsequent sessions.
\end{itemize}

\textbf{Correct speed test:} you can complete every planned repetition without a major pace collapse. If speed drops substantially, foot placement becomes erratic, or you must grab the rails, the work pace is too aggressive.

\subsection{3. Cool-down --- 5 minutes}
Reduce speed gradually to an easy walk. Finish when breathing is controlled and you feel stable.

\section{Workout B: Progressive Threshold Run}

\textbf{Use:} Continuous hard conditioning with less neuromuscular demand than sprint intervals. \textbf{Total time:} 45 minutes.

\begin{center}
\begin{tabularx}{\textwidth}{@{}l l Y@{}}
\toprule
\textbf{Time} & \textbf{Target} & \textbf{Execution}\\
\midrule
0--10 min & RPE 3--4 & Easy progressive warm-up\\
10--20 min & RPE 5--6 & Steady, controlled aerobic running\\
20--30 min & RPE 6--7 & Moderately hard; increase speed or grade slightly\\
30--38 min & RPE 7--8 & Threshold-like effort; controlled discomfort\\
38--40 min & RPE 8--9 & Optional hard finish; \textbf{not} an all-out sprint\\
40--45 min & RPE 2--3 & Easy jog transitioning to walk\\
\bottomrule
\end{tabularx}
\end{center}

Use incline as a load variable rather than a rule. A \textbf{1\%--3\% grade} can increase demand without requiring excessive belt speed, but there is no mandatory incline. Preserve posture and normal stride mechanics.

\section{Progression}

Do not increase speed, incline, repetitions, and work duration simultaneously.

\begin{center}
\begin{tabularx}{\textwidth}{@{}l Y@{}}
\toprule
\textbf{Stage} & \textbf{Progression criterion}\\
\midrule
Weeks 1--2 & Workout A: 6 rounds. Establish repeatable hard speed.\\
Weeks 3--4 & Add a 7th round if all six remain controlled.\\
Weeks 5--6 & Add an 8th round \emph{or} make a small speed/grade increase, not both.\\
Thereafter & Progress only when the previous workload is repeatable with stable mechanics and adequate recovery.\\
\bottomrule
\end{tabularx}
\end{center}

\section{Weekly Placement}

Because hard running creates meaningful musculoskeletal and metabolic load, treat these as \textbf{quality sessions}, not daily conditioning. A strong default is \textbf{1--2 hard treadmill sessions per week}, separated by at least one easier day. If your week also contains boxing, swimming, outdoor running, or demanding lower-body strength work, count the total number of hard lower-body conditioning exposures rather than viewing treadmill HIIT in isolation.

Avoid placing the interval session immediately before another session where fresh legs, footwork, or explosive output are important. Easy treadmill walking can be used on recovery days without turning it into another threshold workout.

\section{Stop / Reduce Intensity Rules}

\begin{tcolorbox}[colback=light,colframe=accent,boxrule=0.8pt,arc=2mm]
Immediately reduce intensity or stop if you develop chest pain or pressure, dizziness, near-fainting, unusual severe breathlessness, loss of coordination, or an inability to control your position on the belt. Use the treadmill safety clip when available.
\end{tcolorbox}

Normal high-intensity discomfort is not the same as losing control. ``Gasping for air'' is not a performance target.

\section{Session Log}

\begin{center}
\begin{tabularx}{\textwidth}{@{}p{0.12\textwidth} p{0.13\textwidth} p{0.13\textwidth} p{0.12\textwidth} Y@{}}
\toprule
\textbf{Date} & \textbf{Workout} & \textbf{Speed/grade} & \textbf{Peak HR} & \textbf{RPE / notes}\\
\midrule
\rule{0pt}{1.7em} & & & &\\
\rule{0pt}{1.7em} & & & &\\
\rule{0pt}{1.7em} & & & &\\
\rule{0pt}{1.7em} & & & &\\
\bottomrule
\end{tabularx}
\end{center}

\section{Technical Notes \& References}

\small
\begin{itemize}
\item Maximum-heart-rate estimate: Tanaka H, Monahan KD, Seals DR. ``Age-predicted maximal heart rate revisited.'' \emph{Journal of the American College of Cardiology}. 2001;37(1):153--156.
\item American College of Sports Medicine. High-Intensity Interval Training educational guidance: high-intensity work intervals commonly use approximately 80--95\% of estimated HRmax.
\item U.S. Centers for Disease Control and Prevention. Measuring Physical Activity Intensity: during vigorous activity, only a few words can generally be spoken without pausing for breath.
\item Heart-rate targets are approximate. Individual response, medication, heat, hydration, fatigue, and measurement method can alter observed heart rate.
\end{itemize}

\vfill
{\footnotesize\color{graytext}
Training handout for exercise programming and self-monitoring; not a diagnostic or medical document. Update the workload from actual performance data rather than age, body mass, or treadmill speed alone.
}

\end{document}
